% Hafta 14 — Her hattan sonra mutlaka iki şey
% Derleme: py -3.12 tools/latex/derle.py docs/week-14/assets/h14-06-test-ve-olcum.tex
\documentclass[tikz,border=8pt]{standalone}
\usepackage{cen429-cizim}
\begin{document}
\begin{tikzpicture}[node distance=8mm and 22mm]
  \node[baslik] (b) at (0,0) {Her hattan sonra mutlaka iki şey};
  \node[altbaslik] at (b.south west) {ölçmeden ``güçlendirdim'' denemez};

  \node[koyukutu=34mm, below=16mm of b.south west, anchor=north west] (n1) {\kb{Gizle}\\Tigress hattı};
  \node[iyikutu=38mm, right=of n1] (n2) {\kb{1· DAVRANIŞI TEST ET}\\birim testleri\\geçmeli};
  \node[uyarikutu=38mm, right=of n2] (n3) {\kb{2· MALİYETİ ÖLÇ}\\boyut · hız\\komut/dal sayısı};
  \node[morkutu=34mm, right=of n3] (n4) {\kb{Karar}\\kabul / geri al};
  \draw[ok, draw=soluk] (n1.east) -- (n2.west);
  \draw[ok, draw=soluk] (n2.east) -- (n3.west);
  \draw[ok, draw=soluk] (n3.east) -- (n4.west);

  \node[dip, text width=170mm, below=8mm of n4.south -| n2, anchor=north] {%
    Örnek ölçüm: \kod{erisim\_ver} 28 komut/3 dal $\rightarrow$ 51 komut/6 dal.};
\end{tikzpicture}
\end{document}
